\subsection{An exact finite cubic bound for the ninth-root unit lattice}

\begin{statement}{A.194}{Proved}
{Four embedded subfield unit groups have elementary quotient of order at most $3^6$}
Retain the embedded fields and roots of A.191.  Let
\(g\in\operatorname{Gal}(M/\Q(z))\) be the actual automorphism
\(g(t)=zt\), and put
\[
 E_{2,j}=g^j(E_2)=F(tz^j),\qquad j=0,1,2.
\tag{A.194.1}
\]
The four cubic intermediate fields of \(M/F\) are
\(F',E_{2,0},E_{2,1},E_{2,2}\).
For a number field \(L\), write
\(\mathcal U_L=\varprojlim_m\mathcal O_L^\times/
 (\mathcal O_L^\times)^{3^m}\).
Define \(B\) to be the closed subgroup of \(\mathcal U_M\)
generated by the images, under these literal field inclusions, of
\(\mathcal U_{F'}\) and the three \(\mathcal U_{E_{2,j}}\).
Then the actual quotient satisfies
\[
        3\mathcal U_M\subseteq B,\qquad
        \mathcal U_M/B\simeq(\Z/3)^b,
                    \qquad 0\leq b\leq6.
\tag{A.194.2}
\]
Here the notation in the first inclusion is additive notation for
the cube map on the completed multiplicative group.
The integer \(b\) is not specified by this bound.
\end{statement}

\begin{proof}
The two generators \(\sigma,h\) in (A.191.2) identify
\(G=\operatorname{Gal}(M/F)\) with \((\Z/3)^2\).
The fixed fields of its four order-three subgroups are exactly
those in (A.194.1) and \(F'\).  For example \(tz^j\)
has cube \(d^2\alpha\zeta^j\), and its stabilizer is
\(\langle\sigma^{-j}h\rangle\).  For these four subgroups
the literal integral group-ring identity is
\[
       \sum_{|H|=3}\ \sum_{a\in H}a
                        =3\cdot1+\sum_{a\in G}a.
\tag{A.194.3}
\]
Each nonidentity element occurs in one order-three subgroup,
while the identity occurs in all four.  Applying this identity
to an actual unit \(u\) gives the exact norm comparison
\[
       \prod_{|H|=3}N_{M/M^H}(u)=u^3N_{M/F}(u)
                         \quad\text{in }\mathcal O_M^\times.
\tag{A.194.4}
\]
Every norm in this formula is followed by its displayed field
inclusion.  The last norm belongs to the base unit group, already
contained in every subfield group.  The maps are continuous for
the pro-three completions, so density gives (A.194.4) there as
well.  It proves \(3\mathcal U_M\subseteq B\).

We specify the finite determinant which gives the numerical
bound.  Use the nine units \(e_1,\ldots,e_9\) of (A.191.7).
Use the A.175 unit conventions
\[
 (u_1,u_2,u_3,u_4,\epsilon,r_2,r_3,r_4)
                                    \quad\text{in }E_2.
\tag{A.194.5}
\]
The pure units \(u_i\) are the literal A.159 units;
\(\epsilon=v/\tau(v)\), with
\(v=(1+\pi)/(1+\pi^2)\), \(\pi=\alpha^2/(\zeta-1)\),
is the A.170 unit.  The \(r_i\) are the specified literal
cube roots of \(\sigma(u_i)/\sigma^2(u_i)\), with
\(r_i\tau(r_i)=1\), in the A.175 coefficient supplement.
These eight units together with \(\zeta\) generate the
full pro-three unit group of \(E_2\).  A.191 proves the
corresponding assertion for the nine \(e_i\).
Since \(\zeta=e_1^3\), it is already supplied by \(F'\).
Consequently the following 33 literal units generate \(B\):
\[
 e_1,\ldots,e_9;\quad
  g^j(u_1),\ldots,g^j(u_4),g^j(\epsilon),
       g^j(r_2),g^j(r_3),g^j(r_4)\quad(j=0,1,2),
\tag{A.194.6}
\]
in precisely this order.

For each row \((p,a,b)\) below, the residue map is the
literal field map \((t,z)\mapsto(a,b)\).
The equations \(a^9=3d^6\), \(b^9=1\), \(b^3\ne1\)
hold in \(\F_p\).  If \(w_i\) is the \(i\)-th unit
of (A.194.6), define its row entry \(m_i\in\{0,1,2\}\)
by
\(w_i(a,b)^{(p-1)/3}=b^{3m_i}\).
All reductions below are nonzero and their displayed
denominators are invertible.
\[
\begin{array}{r|r|r@{\qquad}r|r|r@{\qquad}r|r|r}
 p&a&b&p&a&b&p&a&b\\\hline
307&78&287&307&78&93&307&78&53\\
307&78&168&307&78&274&307&78&46\\
307&98&287&307&98&93&307&143&287\\
307&143&93&307&189&287&523&20&468\\
523&20&410&523&20&217&523&154&468\\
523&154&410&523&156&468&523&156&410\\
523&311&468&523&311&410&613&47&467
\end{array}
\tag{A.194.7}
\]
Read the rows across each line, from left to right.
In the resulting \(21\)-by-\(33\) integer matrix, the
columns
\[
 1,2,3,4,5,6,7,8,9,11,12,15,16,17,19,20,23,24,25,31,32
\tag{A.194.8}
\]
have determinant \(13000436\equiv2\pmod3\).
The certificate evaluates the literal unit coefficients in
these finite fields and verifies this determinant exactly.
Thus the image of \(B\) in \(\mathcal U_M/3\mathcal U_M\)
has dimension at least 21.

Finally \(M\) is totally imaginary of degree 54, so its
free unit rank is 26.  Its three-primary roots of unity
are exactly \(\mu_9\): the actual abelianization
\(\operatorname{Gal}(M/\Q)^{\rm ab}=(\Z/9)^\times\)
in A.191 excludes any larger cyclotomic three-power field.
Hence \(\dim_{\F_3}\mathcal U_M/3\mathcal U_M=27\).
The exact quotient sequence
\[
 0\longrightarrow B/3\mathcal U_M
 \longrightarrow\mathcal U_M/3\mathcal U_M
 \longrightarrow\mathcal U_M/B\longrightarrow0
\tag{A.194.9}
\]
now proves (A.194.2).  This conclusion uses the 21 independent
rows as a lower bound; it does not identify the full image
dimension from a bounded search.

The certificate
\path{code/two_inert_ninth_root_subfield_unit_certificate.py}
verifies the literal
order-nine group-ring identity and all 693 residue entries
of (A.194.7), together with the stated determinant.  The
output and data are retained; the earlier bounded discovery
search is retained separately.
\end{proof}
